\section{A classical shear potential}\label{sm:potential}

Continue each occurrence of $J$ holomorphically from its strict real
argument. Its branch and additive constant are fixed by
\eqref{sm:J}. For a leaf put
\begin{equation}\label{sm:mixed-g}
 g(x,t)=J(2x)-2J(x)-J(x-t)-J(x+t),\qquad
 v=g_t=\Log\frac{\sin\pi(x-t)}{\sin\pi(x+t)}.
\end{equation}
At $t=0$ this is $g(x)$ and $v=0$. The local mixed pyramid has
ordered base
$(1,E/T,-1/T,-E)$, with $E=e^{2\pi\ii x}$ and $T=e^{2\pi\ii t}$.
For small real $t$, its four half-arcs are
$x-t,1/2-x,x+t,1/2-x$. The two positive tetrahedra across its
first diagonal have angle triples
$(x-t,1/2-x,1/2+t)$ and $(1/2-t,x+t,1/2-x)$.
The complementary Lobachevsky terms cancel, giving volume
$\pi g(x,t)$. This fixes the local primitive without a quantum
normalization argument.

For the retained sign $\eps$, define the small hidden parameter
$t=t(x,\chi)$ and its sheared potential by
\begin{equation}\label{sm:shear}
 \chi=t-\frac{\eps v}{2\pi\ii},\qquad
 f_\eps(x,\chi)=g(x,t)-\frac{\eps v^2}{4\pi\ii}.
\end{equation}
At $(a,0)$ the derivative of the first equation in $t$ is
$1-\eps\ii\cot\pi a\ne0$. The holomorphic implicit-function
theorem applies. Differentiating both equations proves
\begin{equation}\label{sm:local-differential}
                         \dd f_\eps=g_x\,\dd x+v\,\dd\chi.
\end{equation}
The construction is separately even in each small $\chi$ because
$g$ is even in $t$ and the shear equation is odd.

Choose independent real coordinates $\mathbf y$ on the original
reduced wire space, keeping every turn and the complementary closing
relation. The affine leaf and gate wires extend to its complexification.
Set
\begin{equation}\label{sm:global-potential}
 \mathcal A_\eps(\mathbf y,\boldsymbol\chi)
       =\sum_i f_{\eps_i}(a_i(\mathbf y),\chi_i)+\sum_{\text{gates}}m.
\end{equation}
At zero color this is exactly $\mathcal F$. Its nondegenerate internal
Hessian is the imported strict variational theorem. A second application
of the holomorphic implicit-function theorem therefore gives a unique
small solution $\mathbf y_*(\boldsymbol\chi)$ of
$\partial_{\mathbf y}\mathcal A_\eps=0$ through the complete optimizer.
Define
\begin{equation}\label{sm:Phi}
 \Phi_\eps(\boldsymbol\chi)
       =\mathcal A_\eps(\mathbf y_*(\boldsymbol\chi),\boldsymbol\chi).
\end{equation}
It is separately even, and the envelope rule gives
$\partial_i\Phi_\eps=v_i$. All these statements concern an explicitly
defined classical germ, independently of any finite-level source.

\begin{lemma}\label{sm:hessian}
With $\tau_i$ the actual primitive ratio in \eqref{sm:tau},
\begin{equation}\label{sm:Hessian}
             \partial_i\partial_j\Phi_\eps(0)
                           =\pi\ii\tau_i\,\delta_{ij}.
\end{equation}
\end{lemma}
\begin{proof}
At a leaf, $g_{tt}(a,0)=-2\pi c$ and
$\chi_t=1-\eps\ii c$. Hence
$f_{\chi\chi}=-2\pi c/(1-\eps\ii c)=\pi\ii\tau$.
Also $g_{xt}(x,0)=0$. Every mixed wire--auxiliary derivative
vanishes at zero, so eliminating the stationary wires adds no
Schur-complement term to this block. Distinct leaves have no direct
mixed auxiliary derivative in \eqref{sm:global-potential}.
\end{proof}

\section{Actual and auxiliary white transitions}\label{sm:maps}

An intermediate mask bit distinguishes two different roles for a
matrix. At a zero bit it is the actual map on the adjacent white
piece. At a one bit there is a shaded face instead; a \emph{completion
map} records the white transition that must be supplied from the other
side when the edge becomes internal. It is not an extra face or a
pairing assigned to that shaded triangle.

\subsection{Universal local formulas}
Normalize the four ordered face flags to $(0,\infty,1,R)$.
For the mixed leaf the relevant coordinates are
\begin{equation}\label{sm:flags}
 R=\frac{\cos^2\pi x}{\cos^2\pi t}
   =\frac{T(1+E)^2}{E(1+T)^2},\qquad
 Z=\frac{1+E}{1+T},\qquad
 Z^2-(1+E)Z+ER=0.
\end{equation}
Its ordered cross ratio is $R^{-1}-1=e^{-g_x}$.
The first apex is $Z$ and the second-copy apex is $Z/E$:
the latter is the reflected minus-$t$ chart.
Reflection here means inversion of the phase coefficients, not
complex conjugation of free parameters. Every defining triple of
face flags is distinct at completeness and remains so nearby.
The two roots $Z,TZ$ of the auxiliary fixed-point equation may
coalesce at completeness; we require no distinctness of that pair
there.

The four white or completion matrices are
\begin{equation}\label{sm:W}
\begin{aligned}
 W_0&=\begin{pmatrix}1&0\\0&E\end{pmatrix},&
 W_1&=\begin{pmatrix}1&E-1\\0&E\end{pmatrix},\\
 W_2&=\begin{pmatrix}E-R&R(1-E)\\E-1&1-ER\end{pmatrix},&
 W_3&=\begin{pmatrix}R&0\\1-E&ER\end{pmatrix}.
\end{aligned}
\end{equation}
All matrices are projective. At a leaf, $W_0,W_2$ are its actual
white maps: they fix the endpoints of the appropriate base edge and
send $Z$ to $Z/E$. The apex equation verifies the statement for
$W_2$. The other two matrices are completion maps at its shaded
cap bases.

We use the rational prism identities of
\cite[equation~(5.5)]{Geometry}, equivalently
\cite[equation~(3.7)]{Fillings}. To specify their charts, put
$\alpha=e^{2\pi\ii(z-x)}$, $\beta=e^{2\pi\ii(x+y-z)}$,
$\gamma=e^{2\pi\ii(z-y)}$. The world prism is
\begin{equation}\label{sm:complex-prism}
\begin{split}
 (v_0,v_1,v_2,v_3)&=(1,\alpha,\alpha\beta,\alpha\beta\gamma),\\
 P&=1+\alpha-
 \frac{\alpha[\alpha\beta(1+\gamma)-1-\alpha]}
      {\alpha^2\beta^2\gamma-\alpha}.
\end{split}
\end{equation}
For an ordered quadruple $(a,b,c,d)$ let
$M(z)=(z-a)(c-b)/((z-b)(c-a))$ and $R=M(d)$, with the
usual limits at infinity. Its world transition is
$(M^\#)^{-1}W_jM$, where $\#$ reciprocates the three phase
variables. Use the ordered tuples \eqref{sm:stack-map}, with
wire phases $\beta\gamma,\alpha\beta,\alpha\beta\gamma$.
The resulting entire maps satisfy
\begin{equation}\label{sm:prism-maps}
\begin{aligned}
 W_{F,1}=W_{G,1}=W_{o,1}&=W_{\rm top},&
 W_{F,3}=W_{G,3}=W_{o,3}&=W_{\rm bottom},\\
 W_{F,2}&=W_{G,0},&
 W_{o,0}=W_{F,0},\qquad W_{o,2}&=W_{G,2}.
\end{aligned}
\end{equation}
The top and bottom maps are the actual white prism triangles.
The cited proof is an identity in independent complex phase variables,
so it applies also to a wire above $1/2$. No completed-child
hypothesis enters that local calculation.

For turns and closure the corresponding identities are
\begin{equation}\label{sm:turn-maps}
\begin{split}
 T_Q(z)&=\frac{R-1}{z-1},\qquad C_Q(z)=\frac{z-R}{z-1},\\
 T_QW_{j+1}(E,R)&=W_j(E^{-1},1-R)T_Q,\\
 C_QW_j(E^{-1},R)&=W_{2-j}(E,R)C_Q .
\end{split}
\end{equation}
They are \cite[equation~(5.6)]{Geometry}. They retain all translations,
not only derivatives at endpoints.

\subsection{The mask induction for the entire maps}
We prove the following strengthened statement simultaneously through
the original expression: at each output, zero bits carry actual white
maps and one bits carry completion maps; at every internal shaded edge
the actual adjacent white piece already has the recorded completion
map. White pieces joined flatly carry the same entire map.

At a child cap seam the prism is white. If the child bit is zero,
\eqref{sm:prism-maps} identifies two actual maps on a flat seam.
If it is one, the shaded triangle remains and the actual prism map
supplies exactly its completion. At the common free edge the
possible bit pairs are $(0,0),(1,0),(0,1)$. In the first case the
actual white maps merge by $W_{F,2}=W_{G,0}$. In the other cases
the unmasked child's actual map supplies the masked child's
completion. The excluded pair $(1,1)$ would require $HH$.
Finally the inherited output free maps and the two actual prism
cap maps give precisely \eqref{sm:mask-product}.

A turn transports this distinction by \eqref{sm:turn-maps}. At the
closing face the strict complementary pair excludes two shaded bits
on the same seam, as already proved. Zero--zero joins actual white
pieces; one--zero supplies the required completion from the actual
white side. No output face remains, so every completion map
that enters a leaf edge return is now its actual neighboring
white transition.

Stationarity of \eqref{sm:global-potential} matches the same complete
complex cross ratios, by \eqref{sm:local-differential}. Three vertices
determine the Möbius map and the cross ratio matches the fourth.
This proves the induction for the deformed charts as well. A deformed
white polygon need not be planar; all its retained triangular pieces
have one common transition by the proved identities.

\section{The original group and the retained positive chain}
\label{sm:group}

\begin{proposition}\label{sm:marked-deformation}
The stationary family \eqref{sm:Phi} determines a holomorphic framed
representation of the original $\pi_1(Y_E^\eps)$ through its complete
holonomy. It retains a compatible positive ideal fundamental chain.
All main peripheral groups remain parabolic, and the original
auxiliary basis has exactly the logarithms \eqref{sm:UV}.
\end{proposition}

We prove the group, chain and peripheral assertions in that order.
In a leaf frame the pure shaded map and the two odd maps are
\begin{equation}\label{sm:shaded}
 S=\begin{pmatrix}0&-ER\\1&-1-E\end{pmatrix},\quad
 A=\begin{pmatrix}R&-R\\1&ER-1-E\end{pmatrix},\quad
 B=\begin{pmatrix}R/E-1-1/E&R\\-1&R\end{pmatrix}.
\end{equation}
They send the original ordered triples as follows:
$S:(Z,\infty,1)\mapsto(Z,0,R)$,
$A:(Z,\infty,1)\mapsto(Z/E,R,0)$, and
$B:(Z,R,0)\mapsto(Z/E,\infty,1)$.
With $S^\#=S(E^{-1},R)$, direct multiplication gives
\begin{equation}\label{sm:edge-words}
\begin{gathered}
 W_0S=S^\#W_0,\quad W_2S=S^\#W_2,\quad W_3S=S^\#W_1,\\
 W_0^{-1}BW_2^{-1}A=I,\quad
 W_2^{-1}BW_0^{-1}A=I,\quad W_1^{-1}BW_3^{-1}A=I .
\end{gathered}
\end{equation}
Equalities are in the projective group. These local identities are
\cite[equations~(5.7)--(5.8)]{Geometry}, for independent $E,R$.
The actual/ completion-map induction is what permits their use on
the enlarged syntax.

The original cut balls have a dual cell decomposition: paired faces
give generators and edge cycles give the relations of its
two-skeleton. Truncating the cusps gives the same group presentation;
higher-dimensional cells add no fundamental-group relation.
Every white seam and every pure or odd shaded return has now been
verified by whole Möbius maps. Hence these maps define an actual
holomorphic framed representation $\rho_{\boldsymbol\chi}$ of
$\pi_1(Y_E^\eps)$ through its oriented complete holonomy. This
conclusion uses the closed topological cut, not a cross-ratio argument
for a hypothetical open four-punctured boundary.

There is also a compatible positive chain on the original manifold.
At completeness first assemble the actual polyhedron $P_E$, whose
ideal vertices are distinct except for the prescribed terminal pairs.
Order these cut vertices and restrict the order to every constituent
cell. Pulling triangulations agree on each shared quadrilateral.
The order selects diagonals; orient every tetrahedron by the ambient
oriented cell, not by its increasing vertex list. Coning from the first
vertex to faces not containing it gives positive nondegenerate
tetrahedra. Give the second projection copy the same labeled
subdivision. White pairings match it, and shaded pairings are already
triangular, so no order on the quotient cusp vertices is required.

Transport these finitely many flags by the actual rational charts.
Their positive complete shapes remain in the positive half-plane
on one small parameter polydisk. The face maps match the subdivided
faces, and edge products stay exactly one. The continuous edge-log
branches have their original sums there. Thus the representation has
the same oriented ideal fundamental chain and boundary triangulation
as the complete metric. We neither require deformed half-balls to
be convex nor triangulate the filled manifold by new positive cells.

\subsection{All main cusps remain parabolic}
At corner $C_j$ the actual small main meridian is
$W_{j-1}^{-1}W_j$ or its inverse. One of these maps may have been
supplied as a completion, but both are now actual white maps.
Their derivatives at the corner agree: they are $E^{-1}$ at
$C_0,C_2$ and $E$ at $C_1,C_3$, using a finite coordinate at
infinity. The meridian multiplier is identically one. The
meridian is nonidentity at completeness and remains so nearby.
It is therefore parabolic. Its commuting peripheral elements
are translations in the same fixed-flag chart. This proves the
statement for every main cusp, including any component closed
inside an intermediate physical product.

\subsection{The exact auxiliary logarithms}
All the based curves \eqref{sm:peripheral-words} remain in the original
leaf collar. Put
\begin{equation}\label{sm:rho-Q}
 \rho=\frac{Z}{1+E-Z}=\frac1T,\qquad
 Q=\frac{E-Z}{Z-1}=\frac{ET-1}{E-T}.
\end{equation}
In a frame sending the selected flag $Z$ to infinity, the
similarity multipliers of the four words are
\begin{equation}\label{sm:multipliers}
        \lambda:Q^2,\qquad \mu_0:\rho,\qquad
        \mu_-:\rho Q,\qquad \mu_+:\rho/Q.
\end{equation}
At a finite fixed flag, this multiplier is the reciprocal of the
Möbius derivative. Substitution of \eqref{sm:flags} into
$\det M/(cZ+d)^2$ verifies these formulas. They are the local
identities \cite[equations~(4.2)--(4.3)]{Fillings}.
Since $\rho=e^{-2\pi\ii t}$ and $Q=e^{-v}$, their small logarithms are
\begin{equation}\label{sm:logs}
 U=-2\pi\ii t+\eps v=-2\pi\ii\chi,\qquad V=-2v.
\end{equation}
These are similarity logs, twice selected eigenvalue logs. Their
constants are zero at completeness. The envelope rule proves
\eqref{sm:UV}; physical turns only conjugate the words.

If $\chi_i=0$, separate evenness gives $V_i=0$ as well. The
auxiliary meridian remains nonidentity nearby, so an auxiliary
cusp left unfilled also remains a genuine parabolic translation
group. All flags and defining face triples stay in a finite
collection of bounded local frames.
This proves Proposition~\ref{sm:marked-deformation}.
